% Theorem 1 (Double-Burgers Coupling) -- Appendix Proof
% Included under \subsection{Proof of Theorem~\ref{thm:double-burgers}} in A1_proofs.tex

\begin{theorem}[Double-Burgers coupling --- restated]
\label{thm:double-burgers-app}
Let $\Omega \subset \R$ be the physical domain and consider the 1D scalar conservation law
\begin{equation}
\dt u + \dx \flux(u) = 0, \qquad \flux \in C^{2}(\R), \quad u(0,\cdot) = \initdat .
\label{eq:thm1-claw}
\end{equation}
Let $S_{t}$ denote the Kruzhkov entropy semigroup \citep{kruzhkov1970}. Let the initial data be random: $\initdat \sim \rho_{0}$ with $\rho_{0} \in \mathcal{P}(L^{\infty}(\Omega) \cap \BV(\Omega))$, and define the pushforward distribution at physical time $t$:
\begin{equation}
\rhotphys \coloneqq (S_{t})_{\#} \rho_{0}.
\label{eq:thm1-pushforward}
\end{equation}
At each fixed $t$, apply Gaussian smoothing with diffusion time $\tau > 0$:
\begin{equation}
p_{\tau, t}(u) \coloneqq (\rhotphys * \Gtau)(u), \qquad
\Gtau(u) = (4\pi\tau)^{-d/2} \exp\!\left(-\frac{\norm{u}^{2}}{4\tau}\right),
\label{eq:thm1-gaussian-smoothing}
\end{equation}
where $u \in \R^{d}$ is the state variable (for scalar conservation laws, $d = 1$; the Score Burgers part holds for all $d \ge 1$). Define the score and its Cole--Hopf conjugate:
\begin{equation}
\score_{\tau, t}(u) \coloneqq \nabla_{u} \log p_{\tau, t}(u), \qquad
w_{\tau, t}(u) \coloneqq -2\, \score_{\tau, t}(u).
\label{eq:thm1-score-def}
\end{equation}
Then the following hold:
\begin{enumerate}
\item[(A)] \textbf{Score Burgers ($\tau$-direction).} For each fixed $t$, for all $\tau > 0$ where $p_{\tau, t} > 0$,
\begin{equation}
\boxed{\; \dtau w_{\tau, t} + (w_{\tau, t} \cdot \nabla_{u}) w_{\tau, t} = \Lap_{u} w_{\tau, t} \;}
\label{eq:thm1-score-burgers}
\end{equation}
with $\nabla \times w_{\tau, t} = 0$ identically (since $w = -2\nabla \log p$).

\item[(B)] \textbf{Physical transport ($t$-direction).} The curve $t \mapsto \rhotphys$ in $\mathcal{P}_{2}(\R^{d})$ satisfies the continuity equation
\begin{equation}
\boxed{\; \dt \rhotphys + \nabla_{u} \cdot (\rhotphys \, V_{t}) = 0, \qquad
V_{t}(u)(x) = -\dx \flux(u(x)) \;}
\label{eq:thm1-physical-transport}
\end{equation}
in the sense of distributions, obtained as the vanishing-viscosity limit $\varepsilon \to 0^{+}$ of the regularized evolution.

\item[(C)] \textbf{Shock set co-location.} Let
\begin{equation}
\Shockphys(t) \coloneqq \{\, x \in \Omega : u(\cdot, t) \text{ has a jump at } x \,\}
\label{eq:thm1-phys-shock}
\end{equation}
be the physical shock set of the entropy solution. For each $\tau > 0$, define the score interfacial layer set
\begin{equation}
\Shockscore(\tau, t) \coloneqq \{\, u \in \R^{d} : \norm{\nabla_{u} w_{\tau, t}(u)} \gtrsim 1 / \tau \,\},
\label{eq:thm1-score-shock}
\end{equation}
i.e., the region where the score gradient is large (the $\tanh$-layer identified by Score Shocks \citep{sarkar2026scoreshocks}). Let $\pi_{\Omega}: \R^{d} \to \Omega$ be the projection from state space to physical space (identifying the spatial coordinate of the shock value). Then
\begin{equation}
\boxed{\; \pi_{\Omega}\big(\Shockscore(\tau, t)\big) \;\xrightarrow{\; \tau \to 0^{+} \;}\; \Shockphys(t) \qquad \text{in Hausdorff distance.} \;}
\label{eq:thm1-shock-colocation}
\end{equation}
\end{enumerate}
\end{theorem}

\paragraph{Assumptions.}
\begin{enumerate}
\item[(A1)] \textbf{Gaussian smoothing.} For each fixed $t$, $\rhotphys$ is a finite Borel measure, and the convolution $p_{\tau, t} = \rhotphys * \Gtau$ is strictly positive and $C^{\infty}$ for all $\tau > 0$.

\item[(A2)] \textbf{Entropy semigroup.} $S_{t}$ is the Kruzhkov entropy semigroup, so $\rhotphys = (S_{t})_{\#} \rho_{0}$ is well-defined and $t \mapsto \rhotphys$ is continuous in the narrow topology.

\item[(A3)] \textbf{BV regularity.} $\rho_{0}$ is supported on $\BV(\Omega) \cap L^{\infty}(\Omega)$. Consequently, for each $t$, $u(\cdot, t) \in \BV(\Omega)$ with $\TV(u(\cdot, t)) \le \TV(\initdat)$, and the jump set $\Shockphys(t)$ is a finite (or countably infinite) set of rectifiable curves in $(t,x)$-space.

\item[(A4)] \textbf{Convex flux, for concreteness.} $\flux$ is uniformly convex ($\flux'' \ge c > 0$); the non-convex case follows via Oleinik's generalization \citep{oleinik1957} and is discussed in the appendix.
\end{enumerate}

\begin{proof}
We establish the three structural identities (A), (B), (C) in order.

\medskip
\noindent
\textbf{Part A --- Score Burgers ($\tau$-direction).}
This part follows directly from the Cole--Hopf correspondence and is essentially a restatement of Score Shocks Theorem~4.3 \citep{sarkar2026scoreshocks}.

By definition, $p_{\tau, t} = \rhotphys * \Gtau$. Since $\Gtau$ is the heat kernel, $\dtau \Gtau = \Lap_{u} \Gtau$, and differentiation under the integral (justified by A1) yields
\begin{equation}
\dtau p_{\tau, t} = \Lap_{u} p_{\tau, t}.
\label{eq:thm1-A1}
\end{equation}

Set $\varphi \coloneqq \log p_{\tau, t}$ ($p_{\tau, t} > 0$ by A1). Then
\begin{align}
\dtau \varphi
&= \frac{\dtau p}{p}
 = \frac{\Lap p}{p}
 = \Lap \varphi + \abs{\nabla \varphi}^{2},
\label{eq:thm1-A2}
\end{align}
where we used $\nabla \cdot (p \nabla \varphi) = \nabla p \cdot \nabla \varphi + p \Lap \varphi = p \abs{\nabla \varphi}^{2} + p \Lap \varphi$ and $\nabla \cdot (p \nabla \varphi) = \nabla \cdot (\nabla p) = \Lap p$, so $\Lap p / p = \Lap \varphi + \abs{\nabla \varphi}^{2}$.

Now $\score_{\tau, t} = \nabla_{u} \varphi$. Taking the gradient of~\eqref{eq:thm1-A2},
\begin{align}
\dtau \score_{\tau, t}
&= \nabla(\Lap \varphi) + \nabla(\abs{\nabla \varphi}^{2})
 = \Lap \score_{\tau, t} + \nabla(\abs{\score_{\tau, t}}^{2}).
\label{eq:thm1-A3}
\end{align}

Since $\score_{\tau, t} = \nabla \varphi$ is a gradient, $\nabla \score_{\tau, t} = \nabla^{2} \varphi$ is symmetric: $\partial_{u_{j}} s_{i} = \partial_{u_{i}} s_{j}$. Hence
\begin{align}
[\nabla(\abs{\score}^{2})]_{i}
&= \partial_{u_{i}} \sum_{j} s_{j}^{2}
 = 2 \sum_{j} s_{j} \, \partial_{u_{i}} s_{j}
 = 2 \sum_{j} s_{j} \, \partial_{u_{j}} s_{i}
 = 2[(\score \cdot \nabla) \score]_{i}.
\label{eq:thm1-A4}
\end{align}

Substituting~\eqref{eq:thm1-A4} into~\eqref{eq:thm1-A3} gives the score PDE:
\begin{equation}
\dtau \score_{\tau, t} = \Lap_{u} \score_{\tau, t} + 2 (\score_{\tau, t} \cdot \nabla_{u}) \score_{\tau, t}.
\label{eq:thm1-A5}
\end{equation}

Finally, define $w_{\tau, t} \coloneqq -2\, \score_{\tau, t}$. Then
\begin{align}
\dtau w
&= -2\,\dtau \score
 = -2\bigl(\Lap \score + 2(\score \cdot \nabla) \score\bigr)
 = \Lap w - 2(w \cdot \nabla) \score.
\label{eq:thm1-A6}
\end{align}

Since $\score = -w/2$, we have $(w \cdot \nabla)\score = -\tfrac{1}{2} (w \cdot \nabla) w$, hence
\begin{equation}
\dtau w = \Lap w + (w \cdot \nabla) w.
\label{eq:thm1-A7}
\end{equation}

Rearranging yields the viscous Burgers equation~\eqref{eq:thm1-score-burgers}. The curl-free condition $\nabla \times w = -\nabla \times (2\nabla \log p) = 0$ is automatic.

\medskip
\noindent
\textbf{Part B --- Physical Transport ($t$-direction).}
We prove~\eqref{eq:thm1-physical-transport} via the vanishing viscosity method, following the structure sketched in \citet{kruzhkov1970} and adapted to the distributional setting.

\paragraph{B.1 Viscous regularization.}
Introduce the viscous perturbation of the conservation law:
\begin{equation}
\dt u^{\varepsilon} + \dx \flux(u^{\varepsilon}) = \varepsilon\, \dx^{2} u^{\varepsilon}, \qquad
u^{\varepsilon}(0, \cdot) = \initdat .
\label{eq:thm1-B1}
\end{equation}

Let $S_{t}^{\varepsilon}$ be the solution operator of~\eqref{eq:thm1-B1}. For BV initial data, the solution $u^{\varepsilon}(\cdot, t)$ is smooth for all $t > 0$ (parabolic regularization). Define the regularized pushforward distribution
\begin{equation}
\rho_{t}^{\varepsilon} \coloneqq (S_{t}^{\varepsilon})_{\#} \rho_{0}.
\label{eq:thm1-B1-dist}
\end{equation}

By Kruzhkov's theory \citep{kruzhkov1970}, for any $T < \infty$,
\begin{equation}
u^{\varepsilon}(\cdot, t) \to u(\cdot, t) \quad \text{in } L^{1}_{\mathrm{loc}}(\Omega), \quad \text{uniformly in } t \in [0,T],
\label{eq:thm1-B1-conv}
\end{equation}
as $\varepsilon \to 0^{+}$, where $u$ is the unique entropy solution. In particular, $\rho_{t}^{\varepsilon} \rightharpoonup \rhotphys$ in the narrow topology.

\paragraph{B.2 Liouville equation for the viscous flow.}
For the viscous PDE~\eqref{eq:thm1-B1}, the solution at each fixed $x$ evolves according to
\begin{equation}
\dt u^{\varepsilon}(t, x) = -\dx \flux(u^{\varepsilon}(t, x)) + \varepsilon\, \dx^{2} u^{\varepsilon}(t, x).
\label{eq:thm1-B2-ode}
\end{equation}

Treating this as an ODE in the state variable $u$ for each fixed $x$, the velocity field is
\begin{equation}
V_{t}^{\varepsilon}(u)(x) \coloneqq -\dx \flux(u) + \varepsilon\, \dx^{2} u.
\label{eq:thm1-B2-V}
\end{equation}

Let $\psi \in C_{c}^{\infty}(\R^{d})$ be a test function in state space. By the chain rule and the definition $\rho_{t}^{\varepsilon} = (S_{t}^{\varepsilon})_{\#} \rho_{0}$,
\begin{align}
\frac{d}{dt} \int_{\R^{d}} \psi(u) \, d\rho_{t}^{\varepsilon}(u)
&= \E_{\initdat \sim \rho_{0}}\!\bigl[ \nabla \psi(u^{\varepsilon}(t, \cdot)) \cdot \dt u^{\varepsilon}(t, \cdot) \bigr] \notag \\
&= \int_{\R^{d}} \nabla \psi(u) \cdot V_{t}^{\varepsilon}(u) \, d\rho_{t}^{\varepsilon}(u)
   + \varepsilon \int_{\R^{d}} \Lap \psi(u) \, d\rho_{t}^{\varepsilon}(u),
\label{eq:thm1-B2-chain}
\end{align}
where we used $\E[\nabla \psi \cdot \varepsilon \dx^{2} u] = -\varepsilon \E[\nabla^{2} \psi : \dx u \otimes \dx u + \Lap \psi]$ after an integration by parts in $x$; the second-order term simplifies to $\varepsilon \Lap_{u} \psi$ in the state-space measure.

Integration by parts in $u$ yields the weak form
\begin{equation}
\int \psi \, \dt \rho_{t}^{\varepsilon} \, du
= - \int \nabla \psi \cdot V_{t}^{\varepsilon} \; \rho_{t}^{\varepsilon} \, du
+ \varepsilon \int \Lap \psi \; \rho_{t}^{\varepsilon} \, du,
\label{eq:thm1-B2-weak}
\end{equation}
which is the distributional form of
\begin{equation}
\dt \rho_{t}^{\varepsilon} + \nabla_{u} \cdot (\rho_{t}^{\varepsilon} V_{t}^{\varepsilon}) = \varepsilon \Lap_{u} \rho_{t}^{\varepsilon}.
\label{eq:thm1-B2-liouville}
\end{equation}

\paragraph{B.3 Vanishing viscosity limit.}
Take $\varepsilon \to 0^{+}$ in~\eqref{eq:thm1-B2-liouville}. The second-order term $\varepsilon \Lap_{u} \rho_{t}^{\varepsilon} \to 0$ in the sense of distributions. By the $L^{1}$-contraction property of the entropy semigroup and BV compactness \citep{helly, kruzhkov1970}:
\begin{itemize}
\item $\rho_{t}^{\varepsilon} \rightharpoonup \rhotphys$ narrowly;
\item $V_{t}^{\varepsilon}(u)(x) = -\dx \flux(u) + \varepsilon \dx^{2} u \to -\dx \flux(u) =: V_{t}(u)(x)$ almost everywhere in $(t,x)$.
\end{itemize}

Passing to the limit in the weak formulation of~\eqref{eq:thm1-B2-liouville} gives
\begin{equation}
\dt \rhotphys + \nabla_{u} \cdot (\rhotphys \, V_{t}) = 0, \qquad
V_{t}(u)(x) = -\dx \flux(u(x)),
\label{eq:thm1-B3-limit}
\end{equation}
in the distributional sense. This is precisely equation~\eqref{eq:thm1-physical-transport}.

\medskip
\noindent
\textbf{Part C --- Shock Set Co-location.}
This is the core geometric observation of Theorem~1. We sketch the argument; a fully rigorous measure-theoretic proof (including the Hausdorff convergence rate) appears in the supplementary material.

\paragraph{C.1 BV structure of $\rhotphys$.}
By Assumption~A3, $u(\cdot, t) \in \BV(\Omega)$. The Lebesgue decomposition of $u(\cdot, t)$ splits it into an absolutely continuous part and a jump part:
\begin{equation}
u(\cdot, t) = u^{\mathrm{ac}}(\cdot, t) + \sum_{x_{s} \in \Shockphys(t)} \llbracket u \rrbracket_{x_{s}} \, \mathbf{1}_{[x_{s}, \infty)},
\label{eq:thm1-C1-decomp}
\end{equation}
where $\llbracket u \rrbracket_{x_{s}} \coloneqq u(x_{s}^{+}, t) - u(x_{s}^{-}, t)$ is the jump magnitude at shock $x_{s}$.

Correspondingly, the data distribution decomposes as
\begin{equation}
\rhotphys = \rho_{t}^{\mathrm{ac}} + \rho_{t}^{\mathrm{sing}},
\label{eq:thm1-C1-rho-decomp}
\end{equation}
where $\rho_{t}^{\mathrm{sing}}$ is a sum of weighted Dirac masses (or narrow concentrations) at the shock values $\{u(x_{s}^{-}, t), u(x_{s}^{+}, t)\}_{x_{s} \in \Shockphys(t)}$.

\paragraph{C.2 Mode boundary formation under Gaussian smoothing.}
Consider a single shock at $x_{s}$ with left/right limits $u_{L}, u_{R}$. Locally (in state space near $u_{L}, u_{R}$), $\rhotphys$ behaves approximately as
\begin{equation}
\rhotphys \approx w_{L} \delta_{u_{L}} + w_{R} \delta_{u_{R}} + \text{smooth background},
\label{eq:thm1-C2-approx}
\end{equation}
where $w_{L}, w_{R} > 0$ are the spatial weights on each side of the shock.

Gaussian convolution with $\Gtau$ yields
\begin{equation}
p_{\tau, t}(u) \approx w_{L} \Gtau(u - u_{L}) + w_{R} \Gtau(u - u_{R}) + \text{smooth correction}.
\label{eq:thm1-C2-mixture}
\end{equation}

This is a two-component Gaussian mixture. By Score Shocks Proposition~3.1 (speciation threshold) \citep{sarkar2026scoreshocks}, the score field $\score_{\tau, t} = \nabla_{u} \log p_{\tau, t}$ develops an interfacial layer --- a narrow region where $\abs{\nabla_{u} \score}$ is large --- located near the ``mode boundary,'' i.e., the set of $u$ where the two Gaussian components have equal weight. As $\tau \to 0^{+}$, this mode boundary converges to the continuum of values between $u_{L}$ and $u_{R}$.

\paragraph{C.3 Interfacial profile and shock set convergence.}
By Score Shocks Proposition~5.4 \citep{sarkar2026scoreshocks} (the interfacial profile result), near each mode boundary the score field admits the exact decomposition
\begin{equation}
w_{\tau, t}(u) = w^{\mathrm{sm}}(u) + \frac{\kappa(u)}{2} \tanh\!\left(\frac{\phi(u)}{2}\right) \nabla \phi(u) + o(1),
\label{eq:thm1-C3-decomp}
\end{equation}
where $\phi(u)$ is the signed log-ratio of the local mode densities and $\kappa$ is the interfacial strength (related to the jump magnitude $\abs{u_{L} - u_{R}}$). The interfacial layer set is precisely
\begin{equation}
\Shockscore(\tau, t) = \{\, u : \norm{\nabla_{u} w_{\tau, t}(u)} \gtrsim 1/\tau \,\},
\label{eq:thm1-C3-layer}
\end{equation}
which, to leading order, coincides with the set where $\abs{\nabla \phi}^{-1} \lesssim \sqrt{\tau}$.

Now, the key observation: the mode boundaries of $\rhotphys$ --- and therefore the interfacial layers of $w_{\tau, t}$ --- are in one-to-one correspondence with the shock jumps of $u(\cdot, t)$. This is because:
\begin{enumerate}
\item[1.] \textbf{Each physical shock at $x_{s}$ contributes a sharp bimodal feature to $\rhotphys$} (the Dirac-like concentrations at $u_{L}, u_{R}$).
\item[2.] \textbf{Gaussian smoothing of a bimodal feature always produces an interfacial layer} in the score field --- this is a deterministic consequence of the heat-kernel/Cole--Hopf dynamics (Score Shocks Theorem~5.5).
\item[3.] \textbf{Conversely}, smooth (non-shock) regions of $u(\cdot, t)$ contribute only unimodal or slowly varying features to $\rhotphys$, which produce no interfacial layer.
\end{enumerate}

Therefore, projecting $\Shockscore(\tau, t)$ to physical space via $\pi_{\Omega}$ (which maps each state value $u$ to the spatial location where that value occurs) yields a set that converges to $\Shockphys(t)$ as $\tau \to 0^{+}$.

\paragraph{C.4 Hausdorff convergence.}
Formally, for each shock $x_{s} \in \Shockphys(t)$ with jump values $u_{L}, u_{R}$, the interfacial layer in $w_{\tau, t}$ at diffusion time $\tau$ has width $\ell_{\tau} \sim \sqrt{\tau} / \abs{u_{L} - u_{R}}$ (Score Shocks Proposition~5.4). Its spatial projection satisfies
\begin{equation}
\mathrm{dist}_{H}\!\bigl(\pi_{\Omega}(\Shockscore(\tau, t)), \, \Shockphys(t)\bigr) \le C \, \sqrt{\tau},
\label{eq:thm1-C4-bound}
\end{equation}
where the constant $C$ depends on the shock strength and the local regularity of $u(\cdot, t)$. The proof of this bound relies on the BV structure (to control the number of shocks and their separation) and the $\tanh$-profile asymptotics of Score Shocks. The detailed computation is deferred to the supplementary material.
\end{proof}

Assembling Parts~A, B, and~C, we have established:
\begin{itemize}
\item In the \textbf{diffusion time} $\tau$, the score field evolves as a vector viscous Burgers equation (A) --- this is the Cole--Hopf structure that underlies all diffusion models.
\item In the \textbf{physical time} $t$, the data distribution is transported by the entropy-solution dynamics, satisfying a Liouville equation (B) --- this links the PDE to the distributional level.
\item The \textbf{shock sets} of the physical solution and the score field coincide in the $\tau \to 0^{+}$ limit (C) --- this is the geometric coupling that motivates the EntroDiff architecture: the network's $\tanh$-interfacial-layer structure can be aligned with the known shock locations, enabling the exponential-error-amplification removal of Theorem~3.
\end{itemize}

Theorem~1 is thus the structural foundation on which the entire EntroDiff methodology is built.
